\begin{helper}[FiberAndDegreeMixedLiftedIntersectionFixedPairInjectionTail]
\label{helper:FiberAndDegreeMixedLiftedIntersectionFixedPairInjectionTail}
Let \(m=\noderef{TwoBiteNaturalReducedVertexCount}(n)\) and
\[
p=\noderef{TwoBiteEdgeProbability}(1/2,n).
\]
Assume \(n\le m^2\) and \(\log n\ge1\).  Fix base graphs \(G_R,G_B\) on
\(\operatorname{Fin}(m)\), a red vertex \(r\), and a blue vertex \(b\).  If
\[
|N_{G_R}(r)|\le2pm,\qquad |N_{G_B}(b)|\le2pm,
\]
then the number of injections
\(\phi:\operatorname{Fin}(n)\hookrightarrow \operatorname{Fin}(m)\times
\operatorname{Fin}(m)\) for which
\[
\left|
\phi(\operatorname{Fin}(n))\cap
\bigl(N_{G_R}(r)\times N_{G_B}(b)\bigr)
\right|
>100(\log n)^3
\]
is at most
\[
\exp(-2(\log n)^3)\,
\left|\{\phi:\operatorname{Fin}(n)\hookrightarrow
\operatorname{Fin}(m)\times\operatorname{Fin}(m)\}\right|.
\]
\end{helper}

\begin{proof}
Write \(L=\log n\), \(T=100L^3\), and
\[
K=N_{G_R}(r)\times N_{G_B}(b)
\subseteq \operatorname{Fin}(m)\times\operatorname{Fin}(m).
\]
The two degree assumptions give
\[
|K|=|N_{G_R}(r)|\,|N_{G_B}(b)|
\le (2pm)(2pm)=4p^2m^2.
\tag{1}
\]
Since \(p=\noderef{TwoBiteEdgeProbability}(1/2,n)\),
\[
p=\frac12\sqrt{\frac Ln}.
\]
The hypothesis \(L\ge1\) implies \(n>0\) and \(L/n\ge0\), so this identity
may be squared, giving \(4p^2n=L\).  Multiplying (1) by \(n/m^2\) therefore
shows that the hypergeometric mean
\[
\mu=n\frac{|K|}{m^2}
\]
satisfies \(\mu\le L\).

By \(\noderef{UniformInjectionImage}\), counting injections whose image has
more than \(T\) points in \(K\) is equivalent, after dividing by the total
number of injections, to counting \(n\)-element subsets of the \(m^2\) product
cells with more than \(T\) points in \(K\).  Choose the standard bijection
between \(\operatorname{Fin}(m)\times\operatorname{Fin}(m)\) and
\(\operatorname{Fin}(m^2)\).  Under this bijection the marked set is a subset
of \(\operatorname{Fin}(m^2)\) of size \(|K|\), and the count just described
is the hypergeometric upper-tail event for population \(m^2\), sample size
\(n\), and marked set \(K\).

Apply
\(\noderef{FiberAndDegreeMixedLiftedIntersectionHypergeometricBound}\) with
population \(m^2\), sample size \(n\), marked set \(K\), threshold \(T\), and
exponential parameter \(1\).  It gives
\[
\Pr(Y\ge T)\le \exp(\mu(e-1)-T).
\tag{2}
\]
Using \(\mu\le L\), \(e-1<2\), and \(L\ge1\), we have
\[
\mu(e-1)-T\le 2L-100L^3\le -2L^3.
\tag{3}
\]
Since the event \(Y>T\) is contained in \(Y\ge T\), (2) and (3) give
\[
\Pr(Y>T)\le \exp(-2L^3).
\]
Multiplying by the total number of injections gives exactly the claimed
cardinality bound.
\end{proof}
